\documentclass{article}

% Language setting
% Replace `english' with e.g. `spanish' to change the document language
\usepackage[deutsch]{babel}

\usepackage{tabularx}

% Set page size and margins
% Replace `letterpaper' with `a4paper' for UK/EU standard size
\usepackage[letterpaper,top=2cm,bottom=2cm,left=3cm,right=3cm,marginparwidth=1.75cm]{geometry}

% Useful packages
\usepackage{amsmath}
\usepackage{graphicx}
\usepackage[colorlinks=true, allcolors=blue]{hyperref}

\title{Pflichtenheft Einführends Projekt THU WS26/27}
\author{Gruppe A Tom Hannig, Lukas Streicher}

\begin{document}
\maketitle

\begin{abstract}
Skillz.
\end{abstract}
\section{Gliederung}
\section{Einleitung}

 Der Zweck dieses Dokuments ist es, eine detaillierte Beschreibung der Anforderungen, sowie der Benutzerschnittstelle für die Anwendung skillz bereitzustellen. 
Es wird abgegrenzt, welche Anforderungen erfüllt werden müssen,damit die entwickelte Anwendung vom Kunden akzeptiert wird. Die App skillz ist eine Anwendung, die es den Benutzern ermöglicht, eine Smartphone App zu benutzen.
Im Folgenden wird auf die Systemfunktionalität und die Systemeinschränkungen, sowie Annahmen über das Produkt eingegangen.
Zuletzt werden detaillierte Anforderungsspezifikation bereitgestellt.

\section{Allgemeine Beschreibung}
Skillz ist eine Smartphone App, die es dem Benutzer ermöglicht Habits zu tracken, Ziele sowie Routinen zu setzen, auf ein persönliches Tagebuch zuzugreifen und endgültig über sich selsbt hinauszuwachsen.
\subsection{Ansichten}

Die App besteht aus mehreren Anischten, sogenannte Screens, über welche der Nutzer mit dem Programm interagieren kann.

\begin{tabularx}{\textwidth}{|l|X|}
        \hline
        \textbf{Screens} & \textbf{Beschreibung} \\
        \hline
        Main-Menu & Nach dem Start der Anwendung befindet sich der Benutzer / die Benutzerin im Main-Menu. Von hier aus kann er / sie seine Skills betrachten filtern und öffnen, die Einstellungen aufrufen, die Goals öffnen, das Journal betrachten, die Routinen aufrufen  Extra Habits tracken, nach Skills/Goals/Routnies suchen oder eine alternative Ansicht des Main-Menus anzeigen lassen.\\
        \hline
        Goals-Screen & Im Goals-Screen kann der Benutzer/die Benutzerin Ziele für sich selbst erstellen, editieren und löschen.Die Goals dienen als langfristige Ziele denen man eine Schwierigkeit zuordnen kann\\
        \hline
        Routines-Screen & Hier kann der Benutzer/die Benutzerin Routinen erstellen,bearbeiten,löschen und vervollständigen.. \\
        \hline
         Journal-Screen & Alle Skills, Extra Skills, Routinen, Goals die an einem Tag eine Veränderung hatten, werden automatisch an diesem Tag ins Journal eingetragen. In diesem Jourunal kann man an dem jeweiligen Tag noch Kommentare hinzufügen und Journal Einträge von vorherigen Tagen lesen. Die Einträge können gefiltert werden und es kann nach bestimmten Einträgen gesucht werden. \\
        \hline
         Extra-Skill-Screen & In diesem screen können besondere Extra Ereingnisse mit einer kurzen Beschreibung hinzugefügt werden.. \\
        \hline
         Einstellungen-Screen & Im Einstellungs-Screen kann der Spieler seine Einstellungen ändern(Ton,Sprache,Account) und ein About Us Pop Up öffnen. \\
   
      
        
    \end{tabularx}

\subsection{Systemeinschränkungen}

Die Anwendung wird durch die Prozessor- und/oder Grafikleistung des Systems begrenzt, auf dem es läuft. 
    Um die Anwendung zu entwickeln, wird Sprache C\# verwendet und in der Version 12.0 benötigt 
    Später soll die Anwendung als Binary für Android und Apple asugeliefert werden.


\section{Spezifische Anforderungen}

Dieser Abschnitt enthält alle spezifsichen Anforderungen an das System. Er bietet eine detaillierte Beschreibung des Systems und seinen Funktionen.



\subsection{Funktionale Anforderungen}

Dieser Abhscnitt enthält alle  Anforderungen, die die grundlegenden Aktion des Softwaresystems spezifieren.


     \begin{tabularx}{\textwidth}{|l|X |} \hline
        \textbf{ID} & \textbf{FA2} \\
        \hline
        TITEL: & Raumschiff anzeigen \\
        \hline
        BESCHREIBUNG: & Das Raumschiff ist ein zentrales Element des Spiels und wird vom Spielenden gesteuert.
        \\
        \hline
        BEGRÜNDUNG: & Das Raumschiff muss für den Spieler / die Spielerin sichtbare angezeigt werden, sodass er / sie sehen kann wie es sich bewegt. \\
        \hline
        ABHÄNGIGKEITEN: & FA1, FA3, FA7, FA21, FA27, FA29, FA37,  FA40, FA42, FA43, FA52, FA53, FA60,FA61, FA66, FA69\\
        \hline
    \end{tabularx}

\subsection{Nichtfunktionale Anforderungen}

  Dieser Abschnitt spezifiziert die nicht-funktionalen Anforderungen an das Softwaresystem.


\section{Zustandsdiagramme}

\LaTeX{} is great at typesetting mathematics. Let $X_1, X_2, \ldots, X_n$ be a sequence of independent and identically distributed random variables with $\text{E}[X_i] = \mu$ and $\text{Var}[X_i] = \sigma^2 < \infty$, and let
\[S_n = \frac{X_1 + X_2 + \cdots + X_n}{n}
      = \frac{1}{n}\sum_{i}^{n} X_i\]
denote their mean. Then as $n$ approaches infinity, the random variables $\sqrt{n}(S_n - \mu)$ converge in distribution to a normal $\mathcal{N}(0, \sigma^2)$.


\section{Klassendiagramm}

Usually the template you're using will have the page margins and paper size set correctly for that use-case. For example, if you're using a journal article template provided by the journal publisher, that template will be formatted according to their requirements. In these cases, it's best not to alter the margins directly.

If however you're using a more general template, such as this one, and would like to alter the margins, a common way to do so is via the geometry package. You can find the geometry package loaded in the preamble at the top of this example file, and if you'd like to learn more about how to adjust the settings, please visit this help article on \href{https://www.overleaf.com/learn/latex/page_size_and_margins}{page size and margins}.


\bibliographystyle{alpha}
\bibliography{sample}

\end{document}
